\section{Transformations generated by flow equations}

Flows in field space build up field transformations
from infinitesimal transformations.
The latter are generally easier to work with than
integral transformations, because they refer to the current
field only. Moreover, the differentiability and invertibility
of the generated transformations is
automatically guaranteed.


\subsection{Flows in field space}

An infinitesimal field transformation,
\begin{equation}
  U\to U+\eps Z(U)U+\rmO(\eps^2),
\end{equation}
is generated by a link field $[Z(U)](x,\mu)$ with
values in $\Lie$. The continuous composition of
such transformations
amounts to integrating a flow equation
\begin{equation}
  \dot{U}_t=Z_t(U_t)U_t
\end{equation}
with respect to a ficticious time $t$. A simple choice
for the generator of the flow is
\begin{align}
  [Z_t(U)]^a(x,\mu)&=\du^a_{x,\mu}\Wl_0(U),
  \\[1.5ex]
  \Wl_0(U)&=\sum_{x,\mu\neq\nu}\,\tr\{U(x,\mu,\nu)\},
\end{align}
where $U(x,\mu,\nu)$ denotes the plaquette loop in the
$(\mu,\nu)$-directions at the point $x$.
In the following, this flow will be referred to
as the ``Wilson flow''.

If $Z_t(U)$ is a differentiable function of $t$ and $U$, the flow equation
(3.2) has a unique solution $U_t$ for any specified initial value $U_0=V$
and all $t\in(-\infty,\infty)$. Moreover,
the solution is differentiable with respect to
$t$ and $V$ (for a proof of these statements,
see ref.~\cite{ArnoldI}, for example).
It should be
emphasized that the existence of the
solution for all times is non-trivial and can
only be guaranteed, without further
assumptions, because the field manifold is compact.


\subsection{Integrated transformations}

At fixed $t$, the field $U_t$ is a well-defined
function of the initial field $V$. Through the
integration of the flow equation, one thus obtains a
differentiable transformation, $V\to U_t=\transt(V)$, of the
field space. The transformation is invertible and its
inverse is differentiable, because
the flow equation can be integrated backwards from
$t$ to $0$.
Moreover,
since the Jacobian $\det\transts(V)$
is equal to unity at $t=0$ and
does not pass through zero at any time,
the transformation is also
orientation-preserving and thus fulfils all
requirements for an acceptable map of field space.

There is a useful compact expression for the Jacobian
which is obtained starting from the equations
\begin{align}
  {\rmd\over\rmd t}\ln\det\transts(V)
  &=\Tr\bigl\{\dot{\trans}_{t,*}(V)\transts(V)^{-1}\bigr\},
  \\[2.0ex]
  [\dot{\trans}_{t,*}(V)](x,\mu;y,\nu)^{ab}
  ={}&-2\,\tr\Bigl\{\hat{\du}^b_{y,\nu}\bigl\{[Z_t(U_t)](x,\mu)U_t(x,\mu)\bigr\}
  U_t(x,\mu)^{-1}T^a
  \notag\\[0.5ex]
  &\hspace{3.2em}-\hat{\du}^b_{y,\nu}U_t(x,\mu)U_t(x,\mu)^{-1}
  [Z_t(U_t)](x,\mu)T^a\Bigr\}.
\end{align}
Noting
\begin{equation}
  \hat{\du}^b_{y,\nu}=\sum_{x,\mu}[\transs(V)](x,\mu;y,\nu)^{ab}
  \du^a_{x,\mu},
\end{equation}
a few lines of algebra then lead to the formula
\begin{equation}
  \ln\det\transts(V)=\int_0^t\rmd s\,\sum_{x,\mu}
  \bigl\{\du^a_{x,\mu}[Z_s(U)]^a(x,\mu)\bigr\}_{U=U_s}.
\end{equation}
In the case of the Wilson flow, for example, the contribution
of the Jacobian to the action of the field $V$,
\begin{equation}
  \ln\det\transts(V)=-\frac{16}{3}\int_0^t\rmd s\,\Wl_0(U_s),
\end{equation}
is proportional to the integral of
the Wilson plaquette action along the flow.
